% Build: `pdflatex main` gives the working draft (system named OpenZTAC);
%        `pdflatex main_anon` (which defines \ANON) gives the double-blind submission.
\ifdefined\ANON
  \documentclass[sigconf,review,anonymous]{acmart}
\else
  \documentclass[sigconf,review]{acmart}
\fi

\settopmatter{printacmref=false}
\setcopyright{none}
\renewcommand\footnotetextcopyrightpermission[1]{}
\acmConference[SACMAT '27]{The 32nd ACM Symposium on Access Control Models and Technologies}{July 19--21, 2027}{Durham, UK}

\usepackage{booktabs}
\usepackage{multirow}
\usepackage{tabularx}
\usepackage{colortbl}
\usepackage{graphicx}
\usepackage{tikz}
\usetikzlibrary{positioning,fit,backgrounds,arrows.meta,calc,shapes.multipart,shapes.geometric,decorations.pathreplacing}
\usepackage[ruled,linesnumbered,vlined]{algorithm2e}
\SetAlFnt{\small}
\SetAlCapFnt{\small}
\usepackage{xspace}
\usepackage{enumitem}
\usepackage{pifont}

\setlength{\emergencystretch}{2em}
\newif\ifanon
\ifdefined\ANON \anontrue \else \anonfalse \fi
% Name of our system: anonymised in the submission build only.
\ifanon
  \newcommand{\sysname}{\textsc{Atrium}\xspace}
\else
  \newcommand{\sysname}{\textsc{OpenZTAC}\xspace}
\fi
\newcommand{\chain}{\textsc{Chain-9}\xspace}
\newcommand{\Pc}{P_{\mathrm{class}}}
\newcommand{\Px}{P_{\mathrm{ctx}}}
\newcommand{\Po}{P_{\mathrm{own}}}
\newcommand{\ow}{\mathsf{ow}}
\newcommand{\Gasg}{G^{\mathrm{asg}}}
\newcommand{\Gdel}{G^{\mathrm{del}}}
\newcommand{\sameDept}{\mathsf{sameDept}}
\newcommand{\sameDiv}{\mathsf{sameDiv}}
\newcommand{\role}{\mathsf{role}}
\newcommand{\decide}{\mathsf{decide}}
\newcommand{\obs}{\mathsf{obs}}
\newcommand{\allowv}{\mathtt{allow}}
\newcommand{\denyv}{\mathtt{deny}}
\newcommand{\emp}{\mathrm{emp}}
\newcommand{\ev}[1]{\mathsf{#1}}
\newcommand{\Routes}{\mathcal{R}}
\newcommand{\ok}{\ding{51}}
\newcommand{\no}{\ding{55}}
\newcommand{\RQ}[1]{\textbf{RQ#1}}
\newcommand{\answer}[2]{\par\smallskip\noindent\fbox{\parbox{\dimexpr\columnwidth-2\fboxsep-2\fboxrule\relax}{%
\textbf{Answer to RQ#1.} #2}}\par\smallskip}
% Notes for the authors: shown in the working draft, removed from the submission build.
\ifanon \newcommand{\TODO}[1]{} \else \newcommand{\TODO}[1]{\textcolor{red}{[#1]}} \fi

\begin{document}

\title[Beyond the Decision Point]{Beyond the Decision Point:\\
Layered Conformance Checking of a Zero-Trust ABAC System\\ (Extended Version)}

\ifanon
\author{Anonymous Author(s)}
\affiliation{\institution{Anonymous Institution}\country{}}
\else
\author{\TODO{Author names}}
\affiliation{\institution{\TODO{Affiliation}}\country{\TODO{Country}}}
\fi

\begin{abstract}
We designed, built and deployed \sysname, a zero-trust system for attribute-based access control (ABAC). Its
decision point applies a chain of nine rules, and the first rule that matches a request decides. We proved
five properties of this chain and tested the decision point thoroughly. Yet when we examined the running
system as a whole, on isolated local deployments, several of these guarantees did not hold. This paper explains why and turns the lesson
into a method. We check a deployment at three levels: the rules (L1), the stored state that the rules read
(L2), and every path that enforces a decision (L3). We prove that, under assumptions we state in one place,
the three levels together guarantee that every effect is authorized by a committed state at most $\Delta$ old,
and that no two of them are enough. At L1, a syntactic condition on the decision code
and a query to a satisfiability modulo theories (SMT) solver make testing on a finite set of cases complete.
At L2, a theorem predicts how long a revoked or expired permission keeps working when grants are cached. At
L3, probes exercise every request path of an inventory generated from the web framework's routing table, and
check its audit record. In \sysname v1, the decision point matched
its model on all 368{,}640 test cases and provably keeps departments apart. The system around it did not
behave as well: the risk rule never took effect, revoked permissions kept working for up to the cache lifetime
(and a race re-admitted them in 44--53\,\% of our trials), denials were never logged,
and several request paths skipped the decision point; four of our eighteen findings appeared only once we
generated the route inventory from the routing table. The repaired \sysname v2 passes all 57 enforcement checks;
in median latency it is 47\,\% slower on document reads and 32\,\% slower on denied file requests, while for
decisions the difference stays within the noise of our runs.
As a preliminary transfer, we also applied parts of the method to OpenFGA (an open implementation of Google's
Zanzibar model), Cedar's TinyTodo and Backstage RBAC. All three agreed with their documented rules on every case
we sampled; in each, the state level exposed a documented behaviour by which a change of permission takes
effect late or not at all, and our theorem predicts OpenFGA's cache window to within 50\,ms.
\end{abstract}

\ccsdesc[500]{Security and privacy~Access control}
\ccsdesc[300]{Security and privacy~Logic and verification}
\ccsdesc[300]{Security and privacy~Software security engineering}

\keywords{attribute-based access control, zero trust, conformance testing, data independence,
differential testing, SMT, authorization caching, revocation, complete mediation, threat model}

\maketitle

% =====================================================================================
\section{Introduction}
% =====================================================================================

Zero Trust Architecture (ZTA) asks a policy decision point (PDP) to evaluate every request against the
current state of subjects, resources and context \cite{rose2020zta,chandramouli2023sp800207a,kerman2025sp180035}.
Attribute-based access control (ABAC) is the usual policy model \cite{hu2014sp800162,jin2012abac}. There is a
long tradition of verifying such policies: change-impact analysis \cite{fisler2005margrave}, encodings of XACML
and cloud policies for satisfiability modulo theories (SMT) solvers \cite{turkmen2017xacmlsmt,backes2018zelkova},
and anomaly detection \cite{hu2011anomaly}. Most recently, Cedar's authorizer was proved correct in Lean and
tested against its production implementation \cite{cutler2024cedar,disselkoen2024vgd}.

All of these results are about a function, $\decide(q)$. Nobody deploys a function. A deployed system reads
attributes and grants from a database, caches them, enforces decisions at several endpoints, and writes an
audit trail. A property of $\decide$ becomes a property of the system only if three more things hold: the
PDP sees the state that the specification talks about; every path to a protected object goes through the PDP,
which is Saltzer and Schroeder's complete mediation \cite{saltzer1975protection}; and the record of each
decision survives. Everyone assumes these. Few check them, and almost never together.

We learned this on our own system. We designed and built \sysname, an open-source ZTA/ABAC system with a
Python backend of about 16.8\,kLOC (FastAPI, PostgreSQL, Redis), a document service with real-time
collaboration, and a decision point that evaluates a first-applicable chain we call \chain: an administrator rule followed by nine
policies ($P_1$, $\Pc$, $P_2$--$P_6$, $\Px$, $P_7$). We
wrote down seven security invariants, proved five properties of the chain on paper, wrote 140 tests that
cover 95\,\% of the PDP, and deployed it. When we finally checked the running system as a whole, the PDP was
exactly right, and the system was not.

\paragraph{What we propose.} This paper describes the method we ended up with and what it found, in
\sysname and elsewhere. We model a deployment by its runs, that is, by the sequences of requests,
decisions, effects, records and commits it can produce (\S\ref{sec:system}), state a threat model over that system (\S\ref{sec:threat}),
and define conformance at three levels (\S\ref{sec:levels}). At \textbf{L1}, the decision function must agree
with its model and satisfy the requirements. At \textbf{L2$_\Delta$}, it must decide on committed state at
most $\Delta$ old. At \textbf{L3}, every effect on a protected object must be preceded by an allow decision
for the right action, its obligations must be met, and the decision must be recorded durably. We write the
policy chain once, as executable code, and interpret it three ways: as Python booleans for differential
testing, as z3 terms for proofs, and as Cedar policies (Figure~\ref{fig:method}).
Put simply, each level asks one question about the deployment:
\begin{itemize}
\item \textbf{L1, rules:} does the decision code match the rule model, and do the rules satisfy the security
requirements?
\item \textbf{L2, state:} does the decision use up-to-date data, or stale data from the cache?
\item \textbf{L3, enforcement:} does every request path actually ask the decision point, and is every decision
logged?
\end{itemize}
For each level we build an \emph{oracle}, a procedure that decides whether the observed behaviour of the system
is correct.

\paragraph{Questions.} We ask five questions. \RQ1: Can rule-level conformance be established
\emph{completely} by bounded checking, and what does it find in \sysname? \RQ2: What goes wrong between stored
state and decisions, and does our staleness bound predict it? \RQ3: Which enforcement failures escape L1 and
L2? \RQ4: Do our existing tests, and each level's oracle, catch faults at that level? \RQ5: Do the fixes the
method suggests work, what do they cost, and how does \sysname, before and after repair, compare with other
authorization systems: OpenFGA, an open implementation of Google's Zanzibar model, Cedar, and Backstage RBAC?
We state all five questions here and answer each in its own section: \RQ1 in \S\ref{sec:l1}, \RQ2 in
\S\ref{sec:l2}, \RQ3 in \S\ref{sec:l3}, \RQ4 in \S\ref{sec:adequacy} and \RQ5 in \S\ref{sec:rq5}.

\paragraph{What is new.} The ingredients have precedents. Cedar proves an authorizer against a model and tests
the implementation against it \cite{cutler2024cedar,disselkoen2024vgd}; model-based security testing checks
that an application enforces a policy \cite{mouelhi2008models}; ACtest checks what a configuration change does to
an application's access decisions \cite{xiang2025actest}; data independence is classical
\cite{wolper1986,lazic2003semantic}. What we add is, first, a \emph{checked} finite abstraction for ordinary PDP
code: a syntactic criterion that a tool verifies on the code and a coverage query that an SMT solver answers,
which together turn a finite differential test of the real decision code into a complete one. Second, a
layered workflow with contracts for state and enforcement, stated over runs, that compose with the rule level
into an end-to-end guarantee and can each be tested on a running deployment. Third, evidence from a system we
built and then repaired in several rounds, each re-checked by the same oracles, that the failures that matter lie where a verified or exhaustively tested PDP
cannot see them. The data-independence argument itself is classical; our contribution there is to make its
premise checkable on ordinary decision code, not a new theory. Our aim also differs from that of
broken-access-control scanners such as BACScan and BACFuzz \cite{liu2025bacscan,dharmaadi2025bacfuzz}, which
find vulnerabilities without a specification: we check a deployment against stated requirements, and report
both what holds and under which assumptions.

\paragraph{Contributions.}
\begin{itemize}
\item A formal model of \sysname as a labelled transition system, a threat model with seven adversary classes
derived from it, and a running example that we follow through every level (\S\ref{sec:system},
\S\ref{sec:threat}).
\item Definitions of L1, L2$_\Delta$ and L3 conformance over runs, with a composition theorem (the three levels
together give end-to-end conformance) and an independence result (no two give the third)
(\S\ref{sec:levels}, \S\ref{sec:discussion}).
\item Two results that make bounded checking sound rather than merely plausible: a data-independence lemma
under a syntactic criterion that a tool checks on the decision code, and a completeness theorem whose premise
an SMT query checks (\S\ref{sec:method}). The coverage check found two holes in the world space we had been
using.
\item A tight staleness theorem for TTL-cached grants (\S\ref{sec:method}), whose prediction the measured windows
match to within 0.63\,s over 100 trials (\S\ref{sec:l2}), and a freshness result for the versioned cache of v2.
\item An L3 check whose route inventory is generated from the web framework's routing table and swept with
subjects who must be refused; it found four failures that our hand-built inventory had missed (\S\ref{sec:l3}).
\item Eighteen findings on \sysname v1, each backed by a proof, an exhaustive count or a reproduction, and
mutation evidence that the faults that matter most are invisible to rule-level oracles and to our own tests
(\S\ref{sec:l1}--\S\ref{sec:adequacy}).
\item \sysname v2, which passes every oracle and keeps all requirements provable at a measured latency cost, an
exact compilation of first-applicable chains to Cedar, and a comparison of \sysname with OpenFGA, Cedar's
TinyTodo, Backstage RBAC and, from the literature, Zanzibar (\S\ref{sec:rq5}).
\end{itemize}
Every experiment ran on isolated local deployments of frozen versions of \sysname; we make no claim about any
other installation. This is one deep case study of a system we built, with preliminary transfer to three
independent systems (\S\ref{sec:compare}). The artifact, with every script, raw result and survivor
classification, is described in Appendix~\ref{app:artifact}.

% =====================================================================================
\section{Background and Related Work}
\label{sec:related}
% =====================================================================================

\paragraph{ABAC and zero trust.} ABAC decides on subject, object and environment attributes
\cite{hu2014sp800162,jin2012abac}; XACML and NGAC are its standards \cite{oasis2013xacml,ferraiolo2016sp800178}.
ZTA adds per-request evaluation by a PDP and enforcement by policy enforcement points (PEPs)
\cite{rose2020zta,chandramouli2023sp800207a,kerman2025sp180035}. Risk-adaptive models feed a risk estimate
into the decision \cite{cheng2007fuzzymls,ameer2022bluesky}. Temporal models and usage control
\cite{bertino2001trbac,park2004ucon} make validity intervals and \emph{ongoing} authorization first-class. The
second is what breaks at our state level.

\paragraph{Analysing policies and authorizers.} Margrave \cite{fisler2005margrave}, SMT-based XACML analysis
\cite{turkmen2017xacmlsmt}, Zelkova \cite{backes2018zelkova} and anomaly detection \cite{hu2011anomaly} reason
about policies. Cedar \cite{cutler2024cedar,disselkoen2024vgd} goes further and proves its authorizer correct,
then tests the Rust code against the Lean model. That is the closest work to our L1. Its guarantees are,
quite reasonably, relative to the entities the application passes in; how fresh those entities are, and
whether every endpoint asks, is left to the integrator. Work on synthesizing policies with LLMs
\cite{vatsa2025synth,vatsa2026autocedar} also targets the policy itself.

\paragraph{Testing access control.} Policy mutation \cite{martin2007mutation} and model-based security testing
\cite{mouelhi2008models} test policies, or applications against policies. RoleCast \cite{son2011rolecast} and
MACE \cite{monshizadeh2014mace} find missing or inconsistent checks without a specification, and ACtest
\cite{xiang2025actest} tests configuration changes in web applications. The three closest pieces of work differ
from ours in what they take and what they give. Cedar's verification-guided development takes a formal model of
the authorizer and its production code, and differential random testing yields confidence that the two agree on
the entities the application passes in; our L1 takes arbitrary PDP code and a rule model, and yields agreement
on every world once criterion G and coverage hold, plus proofs of the requirements, but it says nothing about a
language implementation as Cedar's proofs do. Model-based security testing derives tests of an application from
a security model and qualifies them by mutating the policy; it treats the application's state as given. ACtest
runs a trimmed copy of a production system before and after a configuration change, reports the requests whose
decision flips, and leaves it to a person to say which flips are wanted; it needs no requirements, but it judges
changes, not a deployment. None of the three addresses the freshness of the state a decision reads, which is our
L2. Closest to our L3 are scanners for broken access control. BACScan \cite{liu2025bacscan} crawls a running web
application as several users, links pages that change data to pages that show it, and replays changes as the
wrong user, judging success by the data the application later shows; BACFuzz \cite{dharmaadi2025bacfuzz} fuzzes
PHP applications with parameters chosen by a language model and judges success by the SQL statements they
cause. Both find routes and fixtures automatically and need no specification, which our L3 does not do; in
return, they report vulnerabilities, while our L3 checks each inventoried route against the decision of a
specified PDP and checks liveness and the audit record as well. The two are complementary: a scanner's findings
are candidate routes and fixtures for our inventory. Our L1 oracle builds on differential
and property-based testing \cite{claessen2000quickcheck,maciver2019hypothesis}. Bounded checking is usually
justified by the small-scope hypothesis \cite{jackson2012alloy}; we replace the hypothesis with a
data-independence argument \cite{wolper1986,lazic2003semantic} that we can check.

\paragraph{Authorization state and its freshness.} SAAM and authorization recycling
\cite{crampton2006saam,wei2008recycling} reuse cached or inferred decisions. Zanzibar \cite{pang2019zanzibar}
stores relation tuples, evaluates them through namespace configurations with userset rewrites, and avoids
the ``new enemy'' problem with snapshot reads at the timestamps carried by consistency tokens (zookies).
OpenFGA \cite{openfga} is an open implementation of that model. We bound staleness for the plain TTL cache that
most applications actually use, show which code paths break the bound, evaluate a versioned-key fix in
Zanzibar's spirit, and measure OpenFGA's own check cache against the bound. Table~\ref{tab:related} compares
what each line of work takes and gives; Table~\ref{tab:compare} in \S\ref{sec:compare} compares the systems.

\begin{table*}[t]
\caption{What each line of work takes, what it inspects and what it gives, by conformance level (L1 rule, L2
state, L3 enforcement). ``--'' means the work does not address the column.}
\label{tab:related}
\footnotesize
\begin{tabularx}{\textwidth}{@{}>{\raggedright\arraybackslash}p{3.3cm}@{\;\;}>{\raggedright\arraybackslash}X@{\;\;}>{\raggedright\arraybackslash}X@{\;\;}>{\raggedright\arraybackslash}X@{\;\;}>{\raggedright\arraybackslash}X@{\;\;}>{\raggedright\arraybackslash}X@{}}
\toprule
Work & Takes & Inspects & State freshness (L2) & Paths to objects (L3) & Gives; completeness \\
\midrule
Margrave \cite{fisler2005margrave}, XACML-SMT \cite{turkmen2017xacmlsmt}, Zelkova \cite{backes2018zelkova} & policies & the policy, not the code that runs it & -- & -- & properties of the policy; complete for the policy language \\
Cedar with VGD \cite{cutler2024cedar,disselkoen2024vgd} & a formal model and the production authorizer & the authorizer & entities supplied by the application & -- & proofs about the model; sampled agreement of code and model \\
Policy mutation \cite{martin2007mutation} & a policy and its tests & the policy & -- & -- & adequacy of the tests \\
Model-based security testing \cite{mouelhi2008models} & a security model & the application, through generated tests & taken as given & derived from the model & tests qualified by mutation \\
RoleCast \cite{son2011rolecast}, MACE \cite{monshizadeh2014mace} & application source; no specification & source code, statically & -- & static analysis & candidate missing or inconsistent checks \\
ACtest \cite{xiang2025actest} & a trimmed production copy and a configuration change & the running application & -- & production requests & decisions that flip, for a person to judge \\
BACScan \cite{liu2025bacscan}, BACFuzz \cite{dharmaadi2025bacfuzz} & a running application and test accounts; no specification & HTTP responses (BACScan); SQL statements (BACFuzz) & -- & crawling; fuzzing & vulnerabilities \\
Zanzibar \cite{pang2019zanzibar} & its own design & -- & snapshot reads at token timestamps & clients call \texttt{Check} & a consistency guarantee by design \\
\textbf{This work} & requirements, a rule model, the PDP code and a running deployment & PDP code (criterion G); the running system & bound $\Delta$ proved and measured & routing-table inventory, manual for non-HTTP paths & proofs and counterexamples; L1 complete under G and coverage; L3 relative to the inventory \\
\bottomrule
\end{tabularx}
\end{table*}

% =====================================================================================
\section{\sysname: Design and Formal Model}
\label{sec:system}
% =====================================================================================

\paragraph{Architecture.} Figure~\ref{fig:system} shows the parts of \sysname that matter for authorization.
An administration interface (PAP) manages subjects, objects, grants and the policy chain. The PDP loads
subject and object attributes from PostgreSQL on every decision and reads assignment and delegation grants
through a Redis cache keyed by (subject, object), with lifetimes of 120\,s and 60\,s. Three PEP endpoints serve
files. A document service and a real-time collaboration server have their own entry points. Decisions go to
an append-intended ledger. We analyse version~1 as frozen at the start of the study. \sysname~v2 is the
repaired version that we built from the findings and describe in \S\ref{sec:rq5}. So that each result can be
read in one place, the running example and the result tables and figures of \S\ref{sec:l1}--\S\ref{sec:adequacy}
show v2 next to v1; the reader can ignore the v2 columns until \S\ref{sec:rq5}.

\begin{figure}[t]
\centering
\begin{tikzpicture}[
  font=\scriptsize, node distance=2mm,
  box/.style={draw, rounded corners=1.5pt, align=center, inner sep=2pt, minimum height=5mm},
  pep/.style={box, fill=blue!6, text width=19mm},
  dp/.style={box, fill=green!8, text width=17mm},
  st/.style={box, fill=orange!10, text width=17mm},
  bad/.style={box, fill=red!8, text width=17mm},
  arr/.style={-{Stealth[length=1.5mm]}, semithick}]
  \node[pep] (a1) {\texttt{/access/decide}};
  \node[pep, below=of a1] (a2) {\texttt{/pep/access}\\(+owner bypass)};
  \node[pep, below=of a2] (a3) {\texttt{/pep/preview}, \texttt{/download}};
  \node[pep, below=of a3] (a4) {\texttt{/documents/\{id\}}};
  \node[pep, below=of a4] (a5) {collab.\ WebSocket};
  \node[dp, right=9mm of a2] (pdp) {\chain PDP};
  \node[box, fill=gray!10, text width=17mm, above=3mm of pdp] (pap) {PAP\\(policies, grants)};
  \node[bad, right=9mm of a4] (dr) {document rules};
  \node[bad, right=9mm of a5] (cr) {connect-time check};
  \node[st, right=7mm of pdp, yshift=4mm] (redis) {Redis grant cache\\$\tau_{\mathrm{asg}}{=}120$\,s, $\tau_{\mathrm{del}}{=}60$\,s};
  \node[st, below=4mm of redis] (pg) {PostgreSQL\\subjects, objects, grants};
  \node[st, below=4mm of pg] (log) {decision ledger};
  \draw[arr] (a1) -- (pdp); \draw[arr] (a2) -- (pdp); \draw[arr] (a3.east) -- (pdp.south west);
  \draw[arr] (a4) -- (dr); \draw[arr] (a5) -- (cr);
  \draw[arr] (pdp) -- (redis); \draw[arr] (redis) -- node[right]{miss} (pg); \draw[arr] (pdp) -- (pg);
  \draw[arr] (pap.east) -| (pg.north west);
  \draw[arr, dashed] (pdp.south east) -- (log);
  \draw[arr] (dr) -- (pg); \draw[arr] (cr.east) -- (pg.south west);
  \begin{scope}[on background layer]
    \node[draw=gray, dashed, inner sep=2pt, fit=(pdp)] (L1) {};
    \node[draw=gray, dashed, inner sep=3pt, fit=(redis)(pg)] (L2) {};
    \node[draw=gray, dashed, inner sep=3pt, fit=(a1)(a5)(dr)(cr)(log)] (L3) {};
  \end{scope}
  \node[anchor=south west, gray] at (L1.south east) {L1};
  \node[anchor=south, gray] at (L2.north) {L2};
  \node[anchor=north east, gray] at (L3.south east) {L3};
\end{tikzpicture}
\caption{\sysname~v1: request paths, decision procedures (green: the \chain PDP; red: hand-written rules) and
state. The document API and the collaboration channel never reach the PDP, and the file endpoints skip it for
object owners.}
\label{fig:system}
\end{figure}

\paragraph{Notation.} $\mathbb{T}$ is the set of time instants (non-negative reals, in seconds), $\mathbb{B}=\{\mathit{true},\mathit{false}\}$
the Booleans, $\mathbb{N}$ the natural numbers, and $2^X$ the set of all subsets of a set $X$. $\bot$ marks an
undefined value and is used for nothing else. $[\varphi]$ is 1 if the condition $\varphi$ holds and 0 otherwise.

\paragraph{Entities and state.} Let $U$, $O$, $A$, $D$ and $V$ be the sets of subjects, objects, actions,
departments and divisions, with $\{\mathit{read},\mathit{write},\mathit{download},\mathit{delete}\}\subseteq A$,
and let $\mathit{IP}$ be the set of network addresses. A delegation request has a status in
\[
\mathit{St}=\{\mathit{pending},\mathit{approved},\mathit{rejected},\mathit{expired},\mathit{revoked}\}.
\]
Clearances and sensitivities are totally ordered,
$C=\{\mathit{normal}\prec\mathit{elevated}\prec\mathit{delegated}\}$ and
$S=\{\mathit{public}\prec\mathit{internal}\prec\mathit{sensitive}\prec\mathit{top\_secret}\}$, and are compared
through the ranks $\kappa_C:C\to\{1,2,3\}$ and $\kappa_S:S\to\{0,1,2,3\}$ of the implementation.

\begin{definition}[State]\label{def:state}
A state $\sigma\in\Sigma$ consists of the attribute functions
$\role:U\to\{\mathrm{adm},\mathrm{dir},\mathrm{dm},\emp,\mathrm{cust}\}$, $\mathsf{dept}:U\cup O\to D$,
$\mathsf{div}:U\cup O\to V\cup\{\bot\}$, $\mathsf{clr}:U\to C$, $\mathsf{sens}:O\to S$, $\mathsf{own}:O\to U$, the
risk inputs $\mathsf{fl}:U\to\mathbb{N}$ (failed logins) and $\mathsf{ip}:U\to\mathit{IP}\cup\{\bot\}$ (last known
address), and two grant relations: assignments $\mathit{Asg}\subseteq U\times O\times 2^A\times\mathbb{T}\times
(\mathbb{T}\cup\{\infty\})\times\mathbb{B}$, whose elements $(u,o,B,t_f,t_u,b)$ give the granted action
set $B\subseteq A$, the validity interval $[t_f,t_u)$ and an active flag $b$; and
delegations $\mathit{Del}\subseteq U\times O\times 2^A\times\mathit{St}\times\mathbb{T}$, whose elements
$(u,o,B,s,t_e)$ give the action set, the status and the expiry $t_e$. It further contains the parts that only
procedures other than the PDP read: the policy store $\mathit{Pol}$, which records for each rule whether the PAP
shows it as active; the document workflow, a status
$\mathsf{wf}:O_{\mathrm{doc}}\to\{\mathit{draft},\mathit{review},\mathit{pending},\mathit{approved},
\mathit{published},\mathit{archived}\}$ and collaborators $\mathit{Col}\subseteq U\times O_{\mathrm{doc}}$ for the
objects $O_{\mathrm{doc}}\subseteq O$ that are documents; the open collaboration sessions
$\mathit{Ses}\subseteq U\times O_{\mathrm{doc}}$; and the ledger $\mathit{Led}$, a sequence of decision records.
\end{definition}

The roles are administrator, director, division manager, employee, and $\mathrm{cust}$ for any custom role
that an administrator created. Grants become sets of actions at time $t$:
\begin{align*}
\Gasg_t(u,o) &= \textstyle\bigcup\{B \mid (u,o,B,t_f,t_u,\mathit{true})\in\mathit{Asg},\ t_f\le t<t_u\},\\
\Gdel_t(u,o) &= \textstyle\bigcup\{B \mid (u,o,B,\mathit{approved},t_e)\in\mathit{Del},\ t<t_e\}.
\end{align*}
The delegation workflow always sets $B=\{\mathit{read}\}$. A request is $q=(u,o,a,c)\in Q=U\times O\times A\times
\mathit{Ctx}$, where the context $c\in\mathit{Ctx}$ carries the hour of day $h(c)$ and the source address
$\mathit{ip}(c)\in\mathit{IP}$. The risk score is
\begin{align*}
\rho(u,c)=\min\bigl(100,\;&30\,[h(c)\notin[7,19)]+40\,[\mathsf{fl}(u)\ge5]\\
&+15\,[2\le\mathsf{fl}(u)\le4]+25\,[\bot\neq\mathsf{ip}(u)\neq\mathit{ip}(c)]\bigr),
\end{align*}
and the threshold is $\theta=70$. We write $\sameDept(u,o)$ for
$\mathsf{dept}(u)=\mathsf{dept}(o)$, $\sameDiv(u,o)$ for $\mathsf{div}(u)\neq\bot\wedge\mathsf{div}(u)=\mathsf{div}(o)$,
and $\sameDiv^\circ$ for plain equality of divisions, under which $\bot=\bot$.

\begin{definition}[Rule chain]\label{def:chain}
A chain is a sequence
\[
\Pi=\langle(r_1,e_1,\gamma_1),\dots,(r_n,e_n,\gamma_n)\rangle
\]
of rule names $r_i$, effects $e_i\in\{\allowv,\denyv\}$ and guards $\gamma_i:\Sigma\times\mathbb{T}\times Q\to\mathbb{B}$,
with $\gamma_n\equiv\mathit{true}$ so that some rule always applies.
Its first-applicable semantics is
\[
\decide^\Pi_{\sigma,t}(q)=\bigl(e_k,r_k\bigr)\quad\text{where}\quad k=\min\{i\mid\gamma_i(\sigma,t,q)\}.
\]
\end{definition}

\chain is the chain of Table~\ref{tab:chain}. The administrator rule has no documented priority; it is
implemented as a bypass before the chain, which is the same thing as a first rule.

\begin{table}[t]
\caption{\chain as implemented in v1, first-applicable from top to bottom. $\pi$ is the documented priority.}
\label{tab:chain}
\small
\begin{tabular}{@{}l@{\;}r@{\;}l@{\;}l@{}}
\toprule
Rule & $\pi$ & Eff. & Guard (for $q=(u,o,a,c)$ at time $t$) \\
\midrule
Admin & -- & allow & $\role(u)=\mathrm{adm}$ \\
$P_1$ & 10 & deny & $\neg\sameDept\wedge\Gdel_t(u,o)=\emptyset$ \\
$\Pc$ & 15 & deny & $\kappa_S(\mathsf{sens}(o))>0\wedge\kappa_C(\mathsf{clr}(u))<\kappa_S(\mathsf{sens}(o))$ \\
 & & & $\;\wedge\neg(\mathrm{dir}\wedge\sameDept)\wedge\neg(\mathrm{dm}\wedge\sameDept\wedge\sameDiv^\circ)$ \\
$P_2$ & 20 & allow & $\role(u)=\mathrm{dir}\wedge\sameDept$ \\
$P_3$ & 30 & allow & $\role(u)=\mathrm{dm}\wedge\sameDept\wedge\sameDiv^\circ$ \\
$P_4$ & 40 & allow & $\role(u)\in\{\emp,\mathrm{dm}\}\wedge a\in\Gasg_t(u,o)$ \\
$P_5$ & 45 & deny & $\role(u)=\emp\wedge\Gasg_t(u,o)=\emptyset$ \\
$P_6$ & 50 & allow & $a\in\Gdel_t(u,o)$ \\
$\Px$ & 60 & deny & $\rho(u,c)\ge\theta$ \\
$P_7$ & 99 & deny & $\mathit{true}$ \\
\bottomrule
\end{tabular}
\end{table}

\paragraph{Routes and runs.} The PDP is one of several components that decide. We therefore model the
whole deployment as a labelled transition system: a system state is the committed state together with the
requests in progress, and each transition is labelled by one event. We only need the sequences of labels,
which we call runs. Figure~\ref{fig:run} shows a run and what each conformance level of \S\ref{sec:levels}
will constrain in it.

\begin{definition}[Routes]\label{def:routes}
A route $r\in\Routes$ is an externally reachable operation, given by a decision procedure
$D_r:\Sigma\times\mathbb{T}\times Q\to\{\allowv,\denyv\}$; an action map $\alpha_r:A\cup\{\bot\}\to A$ that gives
the action the procedure evaluates, from the route and possibly from an action $a_c$ sent by the client ($\bot$
if none); a set of obligations $\mathit{Ob}_r$ (such as multi-factor authentication); and an effect: returning
content of $o$ or changing $\sigma$.
\end{definition}

\begin{definition}[Runs]\label{def:runs}
A run is a finite sequence of timestamped events of six kinds, each tagged with the identifier $k$ of the request
it belongs to: $\ev{req}_k(r,u,o,a_c,c)$, a client request on route $r$ by an authenticated principal $u$;
$\ev{dec}_k(u,o,a,d)$, a decision $d$ by some decision procedure;
$\ev{eff}_k(r,u,o,a)$, an effect (content returned, or state changed) on behalf of $u$;
$\ev{rec}(e)$, the durable commit of a record of decision event $e$;
$\ev{com}_k(\Delta\sigma)$, the commit of a state change made by request $k$; and
$\ev{resp}_k$, the response that ends request $k$. Expiry needs no event: it is a function of time. The
committed state $\sigma_t$ is the result of applying, in order, every $\ev{com}$ event with timestamp at most $t$
to an initial state $\sigma_0$. A run is \emph{complete} if every request in it has a response. A run is
\emph{well-formed} if (W1) each request identifier $k$ has exactly one $\ev{req}_k$, which precedes every other
event tagged $k$, and at most one $\ev{resp}_k$, which follows every $\ev{dec}_k$, $\ev{eff}_k$ and $\ev{com}_k$;
(W2) every $\ev{com}_k$ that changes state is preceded, in request $k$, by an $\ev{eff}_k$ for that change, so a
state change, including a change of grants, is an effect that takes hold at its commit; and (W3) every $\ev{rec}(e)$ follows the decision event $e$ it
records. We consider well-formed runs only; the event logs of a web server with one transaction per request,
such as \sysname, have this shape.
\end{definition}

In v1, $D_r$ is the \chain PDP for the decision API and for non-owners on the file endpoints. For owners it
is a constant $\allowv$; for the document API it is a set of hand-written workflow rules; and for the
collaboration channel it is a connect-time check. Figure~\ref{fig:system} draws these four procedures.

\begin{figure}[t]
\centering
\begin{tikzpicture}[font=\scriptsize, >={Stealth[length=1.4mm]}, x=0.86cm,
  ev/.style={circle, draw, inner sep=1.4pt},
  br/.style={decorate, decoration={brace, amplitude=3pt, mirror}, semithick}]
  \node[anchor=east, align=right] at (0,1.1) {admin};
  \node[anchor=east, align=right] at (0,0) {request\\of $u$};
  \draw[gray!60] (0,1.1) -- (7.7,1.1);
  \draw[gray!60] (0,0) -- (7.7,0);
  \node[ev, fill=orange!50] (c1) at (0.7,1.1) {}; \node[above=1pt] at (c1) {$\ev{com}(\Delta\sigma_1)$};
  \node[ev, fill=orange!50] (c2) at (4.9,1.1) {}; \node[above=1pt] at (c2) {$\ev{com}(\Delta\sigma_2)$};
  \node[ev, fill=blue!40] (rq) at (1.2,0) {}; \node[below=1pt] at (rq) {$\ev{req}(r,u,o,a_c,c)$};
  \node[ev, fill=green!50] (dc) at (3.3,0) {}; \node[below=1pt] at (dc) {$\ev{dec}(u,o,a,d)$};
  \node[ev, fill=red!50] (ef) at (5.5,0) {}; \node[below=1pt] at (ef) {$\ev{eff}(r,u,o,a)$};
  \node[ev, fill=gray!50] (rc) at (7.3,0) {}; \node[below=1pt] at (rc) {$\ev{rec}(e)$};
  \draw[->, dashed, green!40!black] (dc) -- (2.3,1.03);
  \node[font=\tiny, green!40!black, anchor=east] at (2.62,0.55) {reads $\sigma_{t'}$};
  \node[font=\tiny, green!40!black, align=left, anchor=west] at (3.2,0.55) {L1: $d=\mathcal{M}(\sigma_{t'},t',q)$\\L2$_\Delta$: $t'\ge t-\Delta$};
  \draw[br] (3.3,-0.42) -- node[below=3pt, font=\tiny]{L3: an allow for $\alpha_r$ and its obligations come first} (5.5,-0.42);
  \draw[br] (3.3,-0.95) -- node[below=3pt, font=\tiny]{L3: the decision is recorded, whatever the outcome} (7.3,-0.95);
  \draw[->, gray] (0,-1.55) -- (7.8,-1.55) node[right]{time};
\end{tikzpicture}
\caption{A run: one request (lower line) interleaved with administrative commits (upper line). The decision
at time $t$ reads the committed state $\sigma_{t'}$. Each level of \S\ref{sec:levels} constrains a different part
of the run: L1 the value of the decision, L2$_\Delta$ how old the state it reads may be, and L3 what must
precede the effect and follow the decision.}
\label{fig:run}
\end{figure}

\subsection{Running Example}
\label{sec:example}

Figure~\ref{fig:example} shows a small organisation that we use throughout. Department \emph{Finance} has the
divisions \emph{Accounting} and \emph{Audit}; department \emph{Legal} has \emph{Compliance}. Ana is an
employee in Accounting with \emph{normal} clearance, Dan manages Audit, and Lea directs Legal. Object $r$ is
an \emph{internal} Accounting report and $o$ an \emph{internal} Legal contract.
\begin{itemize}
\item[(a)] Ana is assigned $\{\mathit{read}\}$ on $r$ and reads it. $P_4$ allows, in v1 and in v2.
\item[(b$_1$)] Dan asks to write $o$. $P_1$ denies, since he is in another department and holds no delegation.
\item[(b$_2$)] Lea now assigns Dan $\{\mathit{read},\mathit{write}\}$ on $o$ and approves his request for a
read-only delegation. In v1 the write passes $P_1$, because a delegation now exists, then $\Pc$, $P_2$ and
$P_3$, and is allowed by $P_4$; \S\ref{sec:l1} shows that this happens in general. In v2, $P_4$ also requires $\sameDept$, the
delegation does not contain \emph{write}, and $P_7$ denies.
\item[(c$_1$)] Ana, unassigned, requests $r$ at 02:00 after five failed logins, so
$\rho=30+40=70\ge\theta$. In v1, $P_5$ denies.
\item[(c$_2$)] Had she been assigned, v1 would have allowed her through $P_4$, and the risk rule would never
have been consulted; \S\ref{sec:l1} proves that this is always so. In v2, $\Px$ is evaluated first and denies both (c$_1$) and
(c$_2$).
\end{itemize}
The v2 column previews the repaired chain $\mathcal{F}$ of \S\ref{sec:rq5}. We computed the traces in
Figure~\ref{fig:example} by evaluating both chains with the reference semantics described in \S\ref{sec:method};
\S\ref{sec:l1} shows that this semantics agrees with the PDP code on every world.

\begin{figure*}[t]
\centering
\begin{minipage}[c]{0.43\textwidth}
\centering
\begin{tikzpicture}[font=\scriptsize, node distance=3mm,
  dept/.style={draw, rounded corners=2pt, fill=blue!8, minimum width=15mm, align=center, inner sep=2pt},
  div/.style={draw, rounded corners=2pt, fill=gray!10, minimum width=15mm, align=center, inner sep=2pt},
  subj/.style={draw, circle, fill=yellow!20, inner sep=0.5pt, minimum size=8mm, align=center},
  obj/.style={draw, fill=orange!12, align=center, inner sep=2pt},
  arr/.style={-{Stealth[length=1.4mm]}, semithick}]
  \node[dept] (fin) at (0,0) {Finance};
  \node[dept] (leg) at (4.6,0) {Legal};
  \node[div] (acc) at (-0.9,-1.0) {Accounting};
  \node[div] (aud) at (1.0,-1.0) {Audit};
  \node[div] (com) at (4.6,-1.0) {Compliance};
  \draw (fin) -- (acc); \draw (fin) -- (aud); \draw (leg) -- (com);
  \node[subj] (ana) at (-1.4,-2.3) {Ana\\$\emp$};
  \node[subj] (dan) at (1.0,-2.3) {Dan\\dm};
  \node[subj] (lea) at (4.0,-2.3) {Lea\\dir};
  \node[obj] (r) at (-0.2,-3.6) {$r$: report\\internal};
  \node[obj] (o) at (3.0,-3.6) {$o$: contract\\internal};
  \draw[gray] (acc) -- (ana); \draw[gray] (aud) -- (dan); \draw[gray] (com) -- (lea);
  \draw[gray] (acc) -- (r); \draw[gray] (com) -- (o);
  \draw[arr, teal] (ana) -- node[left, pos=0.4]{asg \{read\} (a)} (r);
  \draw[arr, purple] (dan) -- node[above, sloped, pos=0.5]{asg \{read, write\}} (o);
  \draw[arr, purple, dashed] (dan.south east) to[bend right=25] node[below, sloped, pos=0.55]{del \{read\}} (o.west);
  \draw[arr, gray, dotted] (lea) -- node[right, pos=0.3]{grants (b$_2$)} (o);
  \node[align=left, anchor=north west, text width=66mm] at (-1.9,-4.3)
    {Ana: \emph{normal} clearance. (c): 02:00, five failed logins, $\rho=70$.};
\end{tikzpicture}
\end{minipage}%
\hfill
\begin{minipage}[c]{0.55\textwidth}
\centering\scriptsize
\input{figs/runex_trace.tex}
\end{minipage}
\caption{Running example. Left: organisation, subjects (role), objects (sensitivity) and grants. Right: rule-by-rule
evaluation of each case in v1 (\chain) and v2 (the repaired chain $\mathcal{F}$); $\cdot$ marks a guard that is
false, the coloured cell the first applicable rule, and blank cells rules that are never reached. Cases:
(a) Ana reads $r$; (b$_1$) Dan writes $o$; (b$_2$) the same after Lea's grants; (c$_1$) Ana, unassigned, reads $r$
under high risk; (c$_2$) the same with an assignment.}
\label{fig:example}
\end{figure*}

\paragraph{Our design specification.} Before this study we had formalized the chain in a design document,
with three differences from the code that define $\mathcal{S}_{\mathrm{doc}}$: (D1) $P_4$ requires
$a=\mathit{read}\wedge\Gasg_t\neq\emptyset$; (D2) $P_3$ and $\Pc$ use $\sameDiv$; (D3) clearance is typed as a
sensitivity, which we resolve with the implementation's ranks. The document proves five properties, among them
an employee least-privilege theorem (LP in Table~\ref{tab:req}).

\paragraph{Requirements.} Table~\ref{tab:req} lists eight requirements, each a predicate $\phi(\sigma,t,q,d)$
over a state, a request and a decision. Since we audit our own system, a reader may reasonably worry that we
chose requirements it fails. To guard against that, we took every requirement word for word from text we had
written before the study: the ``hard invariants'' in our developer guide, the policy descriptions the PAP
displays, and the design document. Appendix~\ref{app:prov} quotes each source. Our own documents disagree in
four places, and as the system's designers we resolved each disagreement ourselves; no separate policy owner
exists for \sysname. We give the reason for each choice and, where it matters, the verdicts under the other
reading. (C1) The invariant ``[employee] access $\leftrightarrow$ assigned or delegated'' gives R4, but $P_5$'s
description denies every unassigned employee. (C2) $\Px$'s description says it ``denies access when the
computed risk score \ldots{} meets or exceeds the configured threshold'' (R6), but the design document places it
after every allow rule. (C3) The developer guide gives the order in which access is decided as ``Admin $\to$
Owner $\to$ Cross-dept+no-delegation $\to$ \ldots'', so an owner may open her own object, yet the invariant on
employee access and $\Pc$'s description name no owner. We follow the statements we had labelled as invariants,
and \S\ref{sec:l1} reports the other readings of C1 and C2. For C3 we keep both statements and read the owner
rule as a narrow exception: an owner may read and download her object, which is what ``opening'' it means,
and only inside her department, because the isolation invariant (R1) has no exception at all. (C4) The
developer guide gives $P_4$ as ``\ldots{} $\wedge$ action=read $\Rightarrow$ allow'' and describes employees as
``read-only on assigned resources within their division'', and the design document says that ``assignments grant read access only,
never write''; $P_4$'s description in the PAP grants access ``only for actions included in the assignment''.
Two documents, one of them the guide, against one: we read assignments as granting read only, which makes the
design document's least-privilege theorem (LP) a requirement, with the owner exception of C3.

\begin{table}[t]
\caption{Requirements, for $\role(u)\neq\mathrm{adm}$, and the documents they come from, all written before the
study (G: developer guide; S: policy descriptions shown by the PAP; M: design document). Appendix~\ref{app:prov}
quotes each source; \S\ref{sec:l1} gives the verdicts. $\ow$ is the owner exception of C3. LP comes from the
design document's theorem, which has no owner exception; with C4 it is a requirement.}
\label{tab:req}
\small
\begin{tabularx}{\columnwidth}{@{}l@{\;}>{\raggedright\arraybackslash}X@{\;}l@{}}
\toprule
 & Requirement & Source \\
\midrule
R1 & $\neg\sameDept\wedge\Gdel_t=\emptyset\Rightarrow\denyv$ (isolation) & G, S \\
R2 & $\mathsf{cls}(u,o)\wedge\neg\ow\Rightarrow\denyv$ (classification) & S, G \\
R3 & $\emp\wedge\allowv\Rightarrow a\in\Gasg_t\cup\Gdel_t\vee\ow$ (need-to-know) & G, S \\
R4 & $\emp\wedge a\in\Gdel_t\wedge\rho<\theta\wedge\neg\mathsf{cls}\Rightarrow\allowv$ (delegation suffices) & G \\
R5 & $\neg\sameDept\wedge\allowv\Rightarrow a\in\Gdel_t$ (cross-dept.\ only via delegation) & G \\
R6 & $\rho\ge\theta\Rightarrow\denyv$ (risk) & S \\
R7 & $\allowv\Rightarrow(\mathrm{dir}\wedge\sameDept)\vee(\mathrm{dm}\wedge\sameDept\wedge\sameDiv)\vee((\emp\vee\mathrm{dm})\wedge a\in\Gasg_t)\vee a\in\Gdel_t\vee\ow$ (every allow within an allow rule's scope) & S, G \\
\midrule
LP & $\emp\wedge\sameDept\wedge\Gdel_t=\emptyset\wedge\allowv\Rightarrow(a=\mathit{read}\wedge\Gasg_t\neq\emptyset)\vee\ow$ (employee least privilege) & M, G \\
\bottomrule
\end{tabularx}
\end{table}

Here $\mathsf{cls}(u,o)$ is $\Pc$'s guard with $\sameDiv$ in place of $\sameDiv^\circ$, and
\[
\ow(u,o,a)\equiv\mathsf{own}(o)=u\wedge\sameDept(u,o)\wedge a\in\{\mathit{read},\mathit{download}\}
\]
is the owner exception of C3. We write $R$ for the set $\{\mathrm{R1},\dots,\mathrm{R7},\mathrm{LP}\}$, and $\mathcal{M}$ for
\chain read as a reference model. The v1 PDP never reads $\mathsf{own}$ (its owner rule sits outside it,
\S\ref{sec:system}), so the exception changes none of its verdicts; we checked this by repeating every L1 run.

% =====================================================================================
\section{Threat Model}
\label{sec:threat}
% =====================================================================================

A zero-trust system is supposed to hold up against its own users. We therefore model an insider, and we
state precisely what the adversary can do and what counts as a successful attack, because the three
conformance levels are defined against exactly these goals.

\paragraph{Assets.} The protected assets are (i) the content of objects, including documents in every
workflow state; (ii) the grant relations $\mathit{Asg}$ and $\mathit{Del}$, since whoever can change them can
change future decisions; and (iii) the decision ledger, which an investigator needs in order to reconstruct who
was allowed and who was refused.

\paragraph{Trusted base.} We trust the administrators (role $\mathrm{adm}$) and the PAP they use, the database,
the cache, the deployed server code and its configuration, and the transport (TLS). Trusting the code means
we assume the server executes what is in the repository; it does not mean we assume the code is correct,
which is what the three levels test. A compromised administrator or database can grant anything, and no
check on the decision logic can detect that.

\begin{definition}[Adversary]\label{def:adv}
The adversary $\mathcal{A}$ controls a set $U_{\mathcal{A}}\subseteq U$ of principals with
$\role(u)\neq\mathrm{adm}$, and holds valid credentials (tokens and sessions) for them. $\mathcal{A}$ has four
capabilities.
\begin{description}
\item[Cap1, requests.] For any $r\in\Routes$, $u\in U_{\mathcal{A}}$, $o\in O$ and $a_c\in A\cup\{\bot\}$,
$\mathcal{A}$ can emit $\ev{req}(r,u,o,a_c,c)$, choosing every field of $c$ that the client controls (body,
headers, query parameters). Server-derived context (time, failed-login count) is not under its control.
\item[Cap2, scheduling.] $\mathcal{A}$ chooses when its requests are sent, including concurrently with any other
principal's requests and transactions, such as an administrator's revocation. It does not control the server's
scheduler.
\item[Cap3, persistence.] Credentials and open sessions issued to $u\in U_{\mathcal{A}}$ stay with $\mathcal{A}$
until they expire. Revoking a grant does not revoke them.
\item[Cap4, workflow.] $\mathcal{A}$ can use every legitimate workflow open to $u$: request delegations, create
and own objects, edit drafts, accept collaboration invitations. Changes that only administrators make, such as
moving $u$ to another department, are not under its control, but it benefits from them.
\end{description}
\end{definition}

\begin{definition}[Attack goals]\label{def:goals}
Let $\Delta^\ast\ge0$ be the staleness that the operator tolerates. A run is an \emph{attack} if it contains
one of the following.
\begin{description}
\item[G1, unjustified effect.] An event $\ev{eff}(r,u,o,a)$ at $t$ with $u\in U_{\mathcal{A}}$ such that
$\mathcal{M}$ denies $(u,o,a,c)$ in $\sigma_{t'}$ for every $t'\in[t-\Delta^\ast,t]$.
\item[G2, requirement violation.] An $\allowv$ for a request $q$ of $u\in U_{\mathcal{A}}$ at $t$ and some
$\phi\in R$ with $\neg\phi(\sigma_{t'},t',q,\allowv)$ for every $t'\in[t-\Delta^\ast,t]$.
\item[G3, unrecorded decision.] A decision event $e$ of request $k$ of some $u\in U_{\mathcal{A}}$ such that
$\ev{resp}_k$ occurs and is not preceded by $\ev{rec}(e)$.
\item[G4, lingering access.] A G1 effect whose only justification in $\mathcal{M}$ was a grant that ended (by
revocation or expiry) before $t-\Delta^\ast$.
\item[G5, unauthorised administration.] A $\ev{com}$ that changes $\mathit{Asg}$ or $\mathit{Del}$ on behalf
of $u\in U_{\mathcal{A}}$ when the administrative policy $\Pi_{\mathrm{adm}}$ (approver, owner, the object's
director, or an administrator) does not authorise $u$.
\end{description}
\end{definition}

G1 and G2 are \emph{historical}: a request is attacked only if no committed state of the last $\Delta^\ast$
justifies it. With caching, no $\Delta^\ast>0$ can promise more, since a grant may end between the state a
decision reads and the decision itself; a current-state guarantee needs synchronous reads. G2 differs from G1
because the model itself may allow the request, and the fault lies in the model. G2 concerns allows only;
R4, which demands an allow, is checked at L1 but is a matter of availability, not an attack goal. G3 matters
because a silently refused prober leaves nothing for an investigator to find.

\paragraph{Out of scope.} Compromise of the trusted base; injection and cross-site attacks on the web
front end; denial of service; side channels; and social engineering. The collaboration server's protocol is in
scope only through its authorization check.

\paragraph{Adversary classes.} Table~\ref{tab:adv} groups the attacks open to $\mathcal{A}$ into seven classes,
according to what the adversary exploits. Each class uses a subset of Cap1--Cap4, reaches one of G1--G5, and has a
\emph{primary} level, the one whose oracle observes it on its own; attacks can combine classes, as a lingering
grant used through a bypass does (Figure~\ref{fig:threat}). We arrived at these classes while analysing v1; the evaluation (\S\ref{sec:eval}) gives a
concrete instance of each, and Table~\ref{tab:findings} at the end of \S\ref{sec:l3} maps every finding to its
class. Figure~\ref{fig:threat} shows where each class enters the system. The mapping already contains our main point:
only A1 is visible to an oracle that tests the decision function.

\begin{table}[t]
\caption{Adversary classes, with the capabilities they use, the goal they reach, and their primary level, whose
oracle observes them on its own.}
\label{tab:adv}
\small
\begin{tabularx}{\columnwidth}{@{}l@{\;}X@{\;}l@{\;}l@{\;}l@{}}
\toprule
 & Class & Cap. & Goal & Lvl \\
\midrule
A1 & over-privileged insider: the chain itself allows too much & Cap1, Cap4 & G2 & L1 \\
A2 & lingering grantee: keeps using a revoked or expired grant & Cap1, Cap3 & G4 & L2 \\
A3 & racer: decides concurrently with a revocation & Cap2 & G4 & L2 \\
A4 & path shopper: picks a route or action the PDP does not see & Cap1 & G1 & L3 \\
A5 & owner and mover: uses ownership, also after leaving the department & Cap1, Cap4 & G1 & L3 \\
A6 & silent prober: probes objects, denials leave no record & Cap1 & G3 & L3 \\
A7 & unauthorised revoker: removes other people's grants & Cap1 & G5 & L3 \\
\bottomrule
\end{tabularx}
\end{table}

\begin{figure}[t]
\centering
\begin{tikzpicture}[font=\scriptsize, node distance=2mm,
  box/.style={draw, rounded corners=1.5pt, align=center, inner sep=1.5pt, minimum height=4.6mm},
  rt/.style={box, fill=blue!6, text width=18mm},
  tb/.style={box, fill=green!8, text width=15mm},
  adv/.style={draw, fill=red!10, rounded corners=2pt, align=center, text width=10mm, minimum height=26mm},
  att/.style={-{Stealth[length=1.4mm]}, red!70!black, semithick},
  lv/.style={font=\tiny\bfseries, fill=white, inner sep=0.6pt}]
  \node[adv] (A) at (0,0) {$\mathcal{A}$\\$U_{\mathcal{A}}$\\(Cap1--Cap4)};
  \node[rt] (r1) at (2.4,1.45) {decision API};
  \node[rt, below=of r1] (r2) {file endpoints};
  \node[rt, below=of r2] (r3) {document API};
  \node[rt, below=of r3] (r4) {collab.\ channel};
  \node[rt, below=of r4] (r5) {grant routes};
  \node[tb] (pdp) at (5.2,0.9) {PDP (\chain)};
  \node[tb, below=3mm of pdp] (cache) {grant cache};
  \node[tb, below=3mm of cache] (db) {database};
  \node[tb, below=3mm of db] (led) {ledger};
  \node[tb, above=3mm of pdp] (adm) {administrators, PAP};
  \begin{scope}[on background layer]
    \node[draw=green!50!black, dashed, thick, rounded corners=3pt, inner sep=2.5pt, fit=(adm)(pdp)(cache)(db)(led), label={[font=\tiny, green!40!black]below:trusted base}] (T) {};
  \end{scope}
  \foreach \i in {1,...,5} \draw[att] (A.east) -- (r\i.west);
  \draw[-{Stealth[length=1.2mm]}, gray] (r1.east) -- (pdp.west);
  \draw[-{Stealth[length=1.2mm]}, gray] (r2.east) -- (pdp.west);
  \draw[-{Stealth[length=1.2mm]}, gray] (pdp) -- (cache); \draw[-{Stealth[length=1.2mm]}, gray] (cache) -- (db);
  \node[lv, text=red!60!black] at ($(r1.east)!0.45!(pdp.west)+(0,0.17)$) {A1 (L1)};
  \node[lv, text=red!60!black] at ($(cache.east)+(0.72,0)$) {A2, A3 (L2)};
  \node[lv, text=red!60!black] at ($(r2.east)+(0.55,-0.16)$) {A4, A5 (L3)};
  \node[lv, text=red!60!black] at ($(r3.east)+(0.42,0)$) {A4 (L3)};
  \node[lv, text=red!60!black] at ($(r4.east)+(0.62,0)$) {A2 (L2), A4 (L3)};
  \node[lv, text=red!60!black] at ($(r5.east)+(0.42,0)$) {A7 (L3)};
  \node[lv, text=red!60!black] at ($(led.east)+(0.5,0)$) {A6 (L3)};
\end{tikzpicture}
\caption{Threat model. The adversary $\mathcal{A}$ uses only its own credentials and the five kinds of route;
everything inside the dashed boundary is trusted. Labels show where each adversary class of
Table~\ref{tab:adv} acts and which level's oracle observes it.}
\label{fig:threat}
\end{figure}

% =====================================================================================
\section{Layered Conformance}
\label{sec:levels}
% =====================================================================================

We now define conformance against the model of \S\ref{sec:system}. A \emph{world} $w=(\sigma,t,q)$ is a state, a
time and a request. Let $\mathcal{I}$ be the decision procedure implemented by the PDP code, read as a function
of the world.

\begin{definition}[Conformance levels]\label{def:levels}
A deployment is
\begin{itemize}
\item \emph{L1-conformant} if for every world $w$, $\mathcal{I}(w)=\mathcal{M}(w)$, for the decision and the
deciding rule, and $\phi(w,\mathcal{M}(w))$ holds for every $\phi\in R$;
\item \emph{L2$_\Delta$-conformant} if in every run, every event $\ev{dec}(u,o,a,d)$ at time $t$ made by the PDP
satisfies $d=\mathcal{I}(\sigma_{t'},t',q)$ for some $t'\in[t-\Delta,t]$; L2 means L2$_0$. We write L2$^+_\Delta$ when
this is required of allow decisions only, and L2$^{c}_\Delta$ (\emph{componentwise}) when the decision equals
$\mathcal{I}$ applied to a \emph{view} that takes the attributes from $\sigma_t$ and each grant set $G^g$ from some
$\sigma_{t_g}$ with $t_g\in[t-\Delta,t]$;
\item \emph{L3-conformant} if in every run, (i) every $\ev{eff}_k(r,u,o,a)$ at $t$ belongs to a request
$\ev{req}_k(r,u,o,a_c,c)$ and is preceded in it by a PDP event $\ev{dec}_k(u,o,a,\allowv)$ for the request's own
principal, object and context, where $a=\alpha_r(a_c)$ and $\alpha_r$ is constant in $a_c$ on every route that
returns content or changes state, and every obligation of that decision is discharged before $t$; (ii) for routes
that change grants, the same holds with $\Pi_{\mathrm{adm}}$ in place of the PDP; and (iii) every $\ev{dec}_k$
event $e$ is followed by $\ev{rec}(e)$ before $\ev{resp}_k$, whatever the outcome of the request.
\end{itemize}
\end{definition}

The first level is a property of a function, the second of a function and its state, the third of the runs.
L3 is a safety property in the sense of Schneider's execution monitoring \cite{schneider2000enforceable}: every
violation of (i)--(iii) is visible in a finite prefix of a run (for (iii), once the response has been sent), so a
monitor, or a probe that observes runs, can detect it.
L2$_\Delta$ implies L2$^+_\Delta$ and L2$^c_\Delta$, but L2$^c_\Delta$ implies neither (\S\ref{sec:method}): a view may
combine grant sets that never held together. The levels compose.

\begin{theorem}[Composition]\label{thm:compose}
If a deployment is L1-, L2$^+_\Delta$- and L3-conformant, then in every run (a) every effect $\ev{eff}(r,u,o,a)$ is
justified: $\mathcal{M}(\sigma_{t'},t',(u,o,\alpha_r,c))=\allowv$ and all $\phi\in R$ hold there, for some $t'$
at most $\Delta$ before the decision that preceded the effect; and (b) every decision is recorded. Hence no run
reaches G1 or G2 with $\Delta^\ast\ge\Delta+\mathit{d}$, where $\mathit{d}$ bounds the time from a decision to its
effect, and no run reaches G3 or G5.
\end{theorem}
\begin{proof}
Take $\ev{eff}_k(r,u,o,a)$ at $t$. By L3(i) it is preceded in request $k$ by $\ev{dec}_k(u,o,a,\allowv)$ at some
$t_d\le t$, with $a=\alpha_r(a_c)$, $q=(u,o,a,c)$ and $t-t_d\le\mathit{d}$. By L2$^+_\Delta$,
$\mathcal{I}(\sigma_{t'},t',q)=\allowv$ for some $t'\in[t_d-\Delta,t_d]$. By L1, $\mathcal{M}$ returns the same
decision in that world, and every $\phi\in R$ holds there. Since $t'\in[t-\Delta-\mathit{d},t]$, the effect is not a
G1 event for $\Delta^\ast\ge\Delta+\mathit{d}$, and the same witness $t'$ shows that the allow is not a G2 event
(every PDP allow is justified in this way, whether or not an effect follows). L3(iii) excludes G3, since a
response is never preceded by a missing record, and L3(ii) excludes G5.
\end{proof}

G4 is a special case of G1, so it is excluded too. The converse question is whether one could test fewer
levels. We return to it in \S\ref{sec:discussion}, once the results have supplied the examples, and show that one
cannot: for each level there is a deployment that conforms at the other two and not at it.


